\subsection{COMPOSITION OPERATORS ON FUNCTIONS OF A REAL VARIABLE}

Let I be an interval in R containing the interval \(R_{+} = [0, +\infty[\) ; let E be a vector space over the field C, formed of functions of a real variable with complex values, defined on I. We suppose that for every \(a \geqslant 0\) and every function \(f \in E\) the function \(x \mapsto f(x + a)\) belongs to E; further, we assume that E contains the restrictions to I of the polynomials with complex coefficients and the exponentials \(e^{\lambda x}\) , where \(\lambda\) is an arbitrary complex number. Any linear map U of E into the space of all maps from I into the field C of complex numbers will be called an operator on E; if \(f \in E\) and \(g = U(f)\) it will be convenient to use the notation

\[
g (x) = U _ {x} ^ {\xi} \left(f (\xi)\right)
\]

where \(\xi\) is a dummy variable in the functional symbolism of the right-hand term (cf.~II, p.~58). For every \(a \geqslant 0\) the operator which to any function \(f \in E\) associates the restriction to I of the function \(x \mapsto f(x + a)\) is called the translation operator by a.

DEFINITION 3. One says that an operator U on E is a composition operator if, for every \(a \geqslant 0\) , it is permutable with the translation operator by a.

In the notation introduced above, this definition becomes the identity

\[
U _ {x + a} ^ {\xi} (f (\xi)) = U _ {x} ^ {\xi} (f (\xi + a))\tag{25}
\]

in \(x\) and \(a\) ( \(x \in \mathrm{I}, a \geqslant 0\) ). One can exchange the rôles of \(x\) and \(a\) in this identity if \(x \geqslant 0\) , then put \(a = 0\) ; one thus obtains for \(x \geqslant 0\)

\[
U _ {x} ^ {\xi} (f (\xi)) = U _ {0} ^ {\xi} (f (\xi + x))\tag{26}
\]

where \(U_{0}\) is the linear form on E which to each function \(f \in E\) associates the value \(g(0)\) of \(g = U(f)\) .

If \(f\) is a polynomial, one has \(f(\xi + x) = \sum_{k=0}^{\infty} \frac{1}{k!} f^{(k)}(\xi)x^k\) , and formula (26) shows that \(U(f)\) is also a polynomial; restricted to the set of polynomials in \(x\) , with coefficients in \(\mathbf{C}\) , (a set which one can identify with the algebra \(\mathbf{C}[\mathbf{X}]\) ), the operator \(U\) is then a composition operator in the sense of def. 1 of VI, p.~270, and all the results of the preceding sections can be applied to it.

We again write \(u_{n}\) for the Appell polynomials attached to the operator U. To the generalized Taylor expansion of a polynomial (VI, p.~274, formula (13)) there corresponds, for more general functions, the following result:

THEOREM 2. Let f be a function admitting a continuous \((n+1)^{th}\) derivative on I, and belonging, with all its derivatives \(f^{(m)}\) for \(1 \leqslant m \leqslant n\) , to E. If U is a composition operator of order \(p \leqslant n\) on E one has, for \(x \geqslant 0\) and \(h \geqslant 0\)

\[
f ^ {(p)} (x + h) = \sum_ {m = 0} ^ {n} \frac {1}{m !} u _ {m} (x) U _ {h} ^ {\xi} \left(f ^ {(m)} (\xi)\right) + \mathrm{R} _ {n} (x, h)\tag{27}
\]

with

\[
\mathrm{R} _ {n} (x, h) = - U _ {h} ^ {\xi} \left(\int_ {0} ^ {\xi - x - h} \frac {1}{n !} u _ {n} (x + \eta) f ^ {(n + 1)} (\xi - \eta) d \eta\right)\tag{28}
\]

(generalized Taylor expansion).

Consider the integral \(\int_{0}^{\xi-x-h}\frac{1}{n!}u_{n}(x+\eta)f^{(n+1)}(\xi-\eta)d\eta\) , defined for all \(\xi\in I\) , and apply to it the formula for integration by parts of order \(n+1\) (II, p.~59, formula (11)); taking account of the relations

\[
u _ {n} ^ {(k)} = n (n - 1) \dots (n - k + 1) u _ {n - k}
\]

derived from (7) (VI, p.~273) by recursion, one obtains

\[
\begin{array}{l} \int_ {0} ^ {\xi - x - h} \frac {1}{n !} u _ {n} (x + \eta) f ^ {(n + 1)} (\xi - \eta) d \eta \\ = \sum_ {m = 0} ^ {n} \frac {1}{m !} u _ {m} (x) f ^ {(m)} (\xi) - \sum_ {m = 0} ^ {n} \frac {1}{m !} u _ {m} (\xi - h) f ^ {(m)} (x + h). \end{array}\tag{29}
\]

We apply the operator U to the two sides of (29), considered as functions of \(\xi\) , then take the value of the function so obtained for the value h of the variable \(\xi\) ; remarking that, by the formulae (26) (VI, p.~277) and (9) (VI, p.~272), one has

\[
U _ {h} ^ {\xi} \big (u _ {m} (\xi - h) \big) = U _ {0} ^ {\xi} \big (u _ {m} (\xi) \big) = \left\{ \begin{array}{l l} 0 & \text {for} m \neq p \\ p! & \text {for} m = p \end{array} \right.
\]

one finally obtains (27).

